\section{Actuadores}

El proceso utiliza los tres tipos de motores requeridos: servomotores para mecanismos de apertura, un motor DC para el cepillado y un motor paso a paso para el transporte del vehículo.

\subsection{Servomotores}
Se emplean dos servomotores. El servo de puerta está conectado al Slave 2 en PD6/D6. La posición programada para puerta cerrada es 145 grados y para puerta abierta es 50 grados. El segundo servo se encuentra en el Slave 1, conectado a PB3/D11, y controla el suministro de líquido; utiliza 0 grados como posición cerrada y 90 grados como posición abierta. Las posiciones son solicitadas por el Maestro mediante comandos enviados por I2C.

\subsection{Motor DC y L298N}
El motor DC destinado al cepillado es controlado por el Slave 1 mediante un puente H L298N. ENA está conectado a PB1/D9, IN1 a PD7/D7 e IN2 a PB0/D8. Durante la fase de lavado, el Maestro solicita dirección hacia adelante y un valor PWM de 128. La fase de cepillado programada dura cinco segundos, tras los cuales el motor se detiene y el vehículo continúa hacia la salida.

\subsection{Motor paso a paso y ULN2003}
El transporte del automóvil se realiza mediante un Stepper 28BYJ-48 controlado por el Slave 2 a través de un ULN2003. Las entradas IN1--IN4 se conectan a PD2, PD3, PD4 y PD5. La librería se inicializa en modo de medio paso y la velocidad utilizada por la secuencia automática es de 800 pasos por segundo.

El Maestro puede ordenar movimiento continuo hacia adelante o en reversa. Durante el proceso normal, el Stepper avanza hasta alcanzar las ventanas de distancia definidas por el VL53L0X, se detiene durante las estaciones de proceso y vuelve a arrancar posteriormente. Al finalizar el ciclo, el programa ejecuta un movimiento en reversa durante cinco segundos antes de esperar el retiro del vehículo.
